\documentclass{article}
\PassOptionsToPackage{numbers,sort&compress}{natbib}
\usepackage[nonanonymous]{neurips_2026}
% Course-proposal adaptation: named authors, no venue notice or review line numbers.
\makeatletter
\renewcommand{\@notice}{}
\makeatother
\AtBeginDocument{\nolinenumbers}
\usepackage[T1]{fontenc}
\usepackage[utf8]{inputenc}
\usepackage{amsmath,amssymb,booktabs,tabularx}
\usepackage{float}
\usepackage{url}
\usepackage{hyperref}
\hypersetup{colorlinks=true,linkcolor=black,citecolor=black,urlcolor=blue,
  pdftitle={Budgeted Molecular Optimization with Reinforcement Learning Fine-Tuning of ChemGPT},
  pdfauthor={Zhang Zhichao; Liao Yang; Han Gaofei}}
\urlstyle{same}
\newcommand{\KL}{D_{\mathrm{KL}}}
\newcommand{\sg}{\operatorname{sg}}

\title{Budgeted Molecular Optimization with\\
Reinforcement Learning Fine-Tuning of ChemGPT}

\author{%
  Zhang Zhichao\qquad Liao Yang\qquad Han Gaofei\\[0.3em]
  {\footnotesize Name order does not imply contribution ranking.}\\[0.4em]
  {\small ME2502 Reinforcement Learning --- Final Project Proposal}
}

\begin{document}
\maketitle

\begin{abstract}
Pretrained chemical language models provide useful molecular priors, but ordinary
sampling does not directly optimize a desired property profile. We propose a
controlled study of whether reinforcement learning improves target-directed
generation under a limited evaluation budget, and how a fixed-prior
Kullback--Leibler (KL) penalty changes the trade-off between optimization and
output quality. Starting from public ChemGPT weights, we will compare frozen
sampling and selection, REINFORCE, and KL-regularized REINFORCE on two GuacaMol
multi-property tasks, using three seeds and matched unique-molecule budgets.
Evaluation will distinguish candidate discovery from the final policy's routine
outputs, while tracking validity, diversity, repetitions, and actual computation.
The project will deliver a reproducible implementation and an analysis of both
improvements and failure modes. In particular, we will examine whether better
search scores correspond to a broadly improved output distribution or to
increasing concentration on a small set of candidates.
\end{abstract}

\section{Background and related work}
Molecular generation proposes structures for de novo design, lead exploration,
and property optimization. Learning a plausible chemical distribution and
finding molecules that satisfy a specified objective are related but distinct
tasks. Useful evaluation therefore includes property quality, structural
coverage, and the number of candidates that must be assessed.

A candidate may be similar to a reference structure yet miss a desired
physicochemical profile, while a high aggregate reward can conceal a weak
component. We therefore consider the joint behavior of the component properties
and the production of multiple distinct candidates. This makes molecular
optimization a useful setting for studying how feedback changes a pretrained
generator under finite evaluation resources.

Existing generators span several representations and mechanisms. Sequence
models learn SMILES or SELFIES autoregressively; ChemGPT extends this approach
through chemical language pretraining~\citep{frey2023scaling}.
Graph and fragment models represent connectivity explicitly: JT-VAE combines
chemical substructures with graph message passing~\citep{jin2018junction}, and
MolGAN uses adversarial graph generation with a property-oriented learning
objective~\citep{decao2018molgan}. Normalizing flows such as MoFlow learn
invertible graph representations~\citep{zang2020moflow}, while equivariant
diffusion models generate atom types and three-dimensional
coordinates~\citep{hoogeboom2022edm}. Representation, generator architecture,
and goal-directed optimization are separate design choices that can be combined.

Goal-directed methods include conditional generation, sampling and filtering,
latent-space optimization, genetic search, and reinforcement learning (RL).
Graph-based genetic algorithms illustrate that relatively simple search can be
competitive on some objectives~\citep{jensen2019ga}. REINVENT established a
chemical-prior-plus-feedback route~\citep{olivecrona2017reinvent};
ChemRLformer further examined language representations and RL design
choices~\citep{ghugare2024chemrlformer}. GuacaMol supplies standardized molecular
design tasks~\citep{brown2019guacamol}, and the Practical Molecular Optimization
(PMO) benchmark emphasizes oracle-call efficiency~\citep{gao2022pmo}.

Our project builds on this established literature by integrating and evaluating
on-policy RL on a fixed public ChemGPT prior. Keeping the generator and scoring
definitions fixed allows us to isolate the effects of parameter updates and
prior regularization. We ask: \textbf{RQ1}, does RL improve discovery relative
to equal-budget frozen sampling and selection? \textbf{RQ2}, does KL
regularization preserve useful output characteristics, and at what cost to the
optimization score? The study will connect policy-gradient learning and prior
regularization to an explicit molecular environment, producing a controlled
comparison that separates the benefits of learning from those of repeated
sampling.

\section{Proposed method}
\paragraph{Model and environment.}
We will use the public \texttt{ncfrey/ChemGPT-4.7M} checkpoint\footnote{\url{https://huggingface.co/ncfrey/ChemGPT-4.7M}.} and its matched
SELFIES/WordPiece tokenizer, without pretraining from scratch. The state is the
current token prefix and time step; an action selects a token or terminates
generation. Each completed molecule receives a deterministic task reward
$R\in[0,1]$. Empty, undecodable, disconnected, or length-truncated outputs receive
zero. The horizon is 256 sampled actions. All methods share the same
checkpoint-compatible start/end tokens, fixed special-token mask, and
temperature-one categorical sampling distribution.

A fixed processing chain will convert sampled tokens into SELFIES, decode them
to SMILES, and apply RDKit validity checks. The resulting canonical identity
will index the score cache and candidate collection. We will retain the original
token sequence, termination status, and property components alongside each
evaluation so that failures and repeated proposals remain visible.

\paragraph{Policy gradient and prior preservation.}
Let $p_\theta(\tau)$ be the probability of a complete sampled token trajectory,
including its termination behavior, and let $p_0$ be the frozen initial policy.
We optimize
\begin{equation}
J_\beta(\theta)=\mathbb{E}_{\tau\sim p_\theta}[R(\tau)]
 -\beta\KL\!\left(p_\theta\Vert p_0\right).
\label{eq:objective}
\end{equation}
Using REINFORCE~\citep{williams1992reinforce}, define $L_i=\log p_\theta(\tau_i)$,
$L_i^0=\log p_0(\tau_i)$, and $G_i=R_i-\beta(L_i-L_i^0)$. For a batch of $N>1$
independent trajectories, the loss is
\begin{equation}
\mathcal{L}(\theta)=-\frac{1}{N}\sum_{i=1}^{N}\sg(G_i-b_i)L_i,
\qquad b_i=\frac{1}{N-1}\sum_{j\ne i}G_j,
\label{eq:loss}
\end{equation}
where $\sg$ denotes stop-gradient. The sampled log-ratio accounts for the
trajectory KL in Eq.~\eqref{eq:objective}; no additional KL loss is added.
A singleton tail batch uses a previous-batch baseline. We will log full-vocabulary
prefix KL as a diagnostic, keeping the training objective unchanged. This regularizer acts on token
trajectories rather than directly on the set of canonical molecules. Different
encodings can map to the same molecule, so their sampled probabilities must
remain attached to their original trajectories. The baseline uses other samples
in the batch to reduce variance without depending on the current sample's action.

Each fresh on-policy batch receives one update. Repeated molecules retain their
sampled trajectories and cached rewards in the loss; canonicalization is used
for scoring identities, not to replace training trajectories. Initial settings
are batch size 32, Adam learning rate $10^{-5}$, gradient norm clipping at 1.0,
and $\beta=0.01$ for the regularized model. Actor and reference likelihoods use
the same token mask, temperature, and trajectory convention; the reference
parameters remain frozen. Development checks will precede a configuration freeze.
We will verify that padding leaves trajectory probabilities unchanged, the
reference receives no gradients, and rollout probabilities match their
training-time recomputation. These checks will be completed before running the
main comparisons.

\clearpage
\section{Experimental design}
\paragraph{Tasks and controls.}
We will use GuacaMol's goal-directed \textbf{v2} definitions, whose selected MPO
rewards use geometric aggregation. Table~\ref{tab:tasks} separates development
from the two main comparisons. Scores combine reference similarity and
calculated physicochemical properties; they are proxies for design requirements,
not measured binding or efficacy.

\begin{table}[H]
\caption{Public tasks and their roles. MPO denotes multi-property optimization.}
\label{tab:tasks}
\centering
\begin{tabularx}{\linewidth}{@{}l l X@{}}
\toprule
Task & Role & Main scoring components \\
\midrule
Perindopril MPO & Development & Reference similarity and aromatic ring count \\
Fexofenadine MPO & Main comparison & Atom-pair similarity, TPSA, and logP \\
Osimertinib MPO & Main comparison & Two fingerprint similarities, TPSA, and logP \\
\bottomrule
\end{tabularx}
\end{table}

Fexofenadine MPO favors reference similarity and sufficiently high TPSA while
limiting logP. Osimertinib MPO combines similarity requirements from two
fingerprints with a different physicochemical profile. Their components will be
scored using the published definitions, with no task-specific rescaling from
main-run outcomes. Each task will be optimized independently from the same
initial checkpoint.

The three conditions are \textbf{Frozen + best-of-$B$}, retaining the best
observed candidates without parameter updates; \textbf{REINFORCE}, using
$\beta=0$; and \textbf{REINFORCE + KL}, using the fixed-prior penalty. Two tasks,
three conditions, and three seeds yield \textbf{18 runs}. Development will use
Perindopril and separate seeds. Fexofenadine has already been explored in an
earlier, different GRU pilot; this history will be disclosed rather than
presented as a fresh blind test. Main-task outcomes will not select
hyperparameters, seeds, or checkpoints. For each search curve, frozen and
trainable methods will be compared using the same top-$k$ statistic over
candidates available at that point. The frozen baseline will receive the full
evaluation budget under the same sampling protocol.

\paragraph{Evaluation budget.}
Each run may score up to $B=3{,}000$ previously unseen valid canonical molecules,
with additional caps of 15,000 generation attempts, 470 optimizer updates, and
30 minutes of search/training time. The first exhausted limit ends the run.
The tail batch size is determined before sampling by remaining unique-query
capacity. Invalid and repeated attempts are recorded; caching never grants
another run free logical queries. Unique-query matching does not imply equal
training computation, so updates, tokens, physical scoring calls, and wall time
will also be reported. The evaluation ledger will distinguish generation
attempts, first-time valid identities, cache hits, and actual scoring calls.
This will reveal whether a method spends substantial computation revisiting
known candidates even when its unique-query count appears efficient.

\paragraph{Outcomes and analysis.}
The primary endpoint is the normalized area under the top-10 MPO score curve
against new unique evaluations. We will also report budget-end scores, the
number of threshold-passing candidates, and strict joint component satisfaction;
an aggregate MPO threshold does not imply that every component passes. For
early-stopped runs, fixed-budget summaries will explicitly identify any
last-value carry-forward and show the achieved query count.

Each terminal policy will independently generate 1,024 additional attempts,
without further training or checkpoint selection. We will measure all-attempt
mean reward, validity, unique valid output, repetition, and fingerprint/scaffold
diversity. Search histories and final-policy samples will be analyzed separately.
Results will show every seed, means and variability, and paired differences for
each task. A higher search score accompanied by collapsed routine outputs, or a
KL model trading peak score for diversity, will be analyzed as a trade-off.
We will inspect component-score and sequence-length distributions alongside KL
and repetition curves to explain these outcomes. Planned figures will include
budget--quality curves, paired method comparisons, and a summary of the
quality--diversity trade-off.

\clearpage
\section{Feasibility, resources, and work plan}
\paragraph{Feasibility and resources.}
Initial CPU checks loaded the public ChemGPT checkpoint, generated 32 naturally
terminated and decodable single-molecule samples, verified rollout/teacher-forcing
log-probability agreement, and obtained finite nonzero gradients. An enumerable
two-step environment checked the KL score-function and leave-one-out gradient
construction. H100 resource profiling has also been completed. These are
feasibility checks; the formal 18-run comparison remains planned. The next
development milestone is an end-to-end short run that validates the scorer,
query ledger, and frozen-reference behavior together. A small development
configuration will then be fixed and reused across both main tasks and all
reported seeds.
The planned computing resources are one NVIDIA H100 80GB GPU, 4 vCPUs,
64 GiB of host RAM, and 10 GB of project storage.

\paragraph{Team responsibilities.}
\textbf{Zhang Zhichao} will cover molecular-generation literature, research-question
formulation, and overall study design; adapt ChemGPT loading, sampling, and
decoding interfaces for integration with the RL framework; and coordinate the
proposal's background, motivation, and overall narrative in the final report.
This work will provide the tested model interface and the study's scientific
framing.

\textbf{Liao Yang} will develop the GuacaMol task, property-scoring, and reward
interfaces; implement evaluation budgets, duplicate-molecule caching, and the
frozen-model baseline; organize comparison runs; and carry out metric
verification, result analysis, figures, and experimental reporting. The
corresponding outputs will be an auditable scoring ledger and reproducible
comparison tables.

\textbf{Han Gaofei} will focus on RL-method design and implementation, including
REINFORCE, the variance-reduction baseline, the frozen reference policy, and KL
regularization; verify policy updates, gradient correctness, and training
stability; develop training configurations; and write the methods section and
interpret the optimization behavior. This work will deliver the validated
policy-gradient training module and its implementation checks.

\paragraph{Joint work and deliverables.}
All three members will jointly finalize experiments, cross-check code and
results, discuss conclusions, integrate the report, prepare the presentation,
and participate in Q\&A. Deliverables are versioned code and configurations,
complete generation/scoring logs, per-seed metrics, necessary checkpoints,
and a reproducible final report. Negative results and quality--diversity
trade-offs will be retained. Each implementation component will be reviewed by
another member, and the team will jointly rehearse explanations of the MDP,
objective, and empirical findings. The schedule is summarized in
Table~\ref{tab:schedule}.

\begin{table}[H]
\caption{Planned schedule and concrete outputs.}
\label{tab:schedule}
\centering
\begin{tabularx}{\linewidth}{@{}l X@{}}
\toprule
Period & Output \\
\midrule
Before November 8 & Proposal, development checks, and frozen configuration \\
November 9--15 & Complete the 18-run core comparison \\
November 16--22 & Independent metric checks and one focused sensitivity study \\
By December 6 & Final report, reproducibility package, and presentation materials \\
\bottomrule
\end{tabularx}
\end{table}

\clearpage
\bibliographystyle{unsrtnat}
\bibliography{references}
\end{document}
